\chapter{As leis da corrente e da tensão nos circuitos}\label{ch:g8:circuit-laws}

Dois aparelhos, duas grandezas --- e agora, enfim, a constituição.
Todo circuito que você vier a montar, da lanterna à rede de uma
cidade, obedece a quatro leis curtas: duas sobre a contabilidade da
corrente, duas sobre a da \omterm{def:g8:voltage-voltmeter:voltage}{tensão}. Armado com elas e um pouco de
álgebra, você consegue calcular leituras que ainda não mediu --- e
pegar qualquer medida que minta.

\section{As leis da corrente}

\begin{proposition}[Lei da corrente numa volta em série]\label{prop:g8:circuit-laws:seriesI}
Ao longo de uma única volta, a intensidade é a mesma em todo lugar:
\[
I_{\text{pilha}} = I_{\text{lâmpada 1}} = I_{\text{lâmpada 2}} =
\dots
\]
Uma estrada, uma marcha --- os pedágios do capítulo passado, agora
uma lei da constituição.
\end{proposition}

\begin{proposition}[A lei dos nós]\label{prop:g8:circuit-laws:junction}
Em qualquer \omterm{def:g5:series-parallel-circuits:parallel}{nó}, as intensidades que chegam juntas são iguais às
intensidades que saem juntas: para uma estrada principal que se
divide em dois \omterm{def:g5:series-parallel-circuits:parallel}{ramos},
\[
I = I_1 + I_2 .
\]
Carga nenhuma é criada, destruída ou estocada numa bifurcação: tudo o
que chega, sai. (Vale para qualquer número de \omterm{def:g5:series-parallel-circuits:parallel}{ramos} --- some as
chegadas, some as partidas, e as contas fecham.)
\end{proposition}

\begin{omfigure}
\begin{tikzpicture}[scale=0.85]
  \draw (0,0) to[battery1] (0,3.2)
    to[ammeter, l={$I = \qty{0.50}{A}$}] (3.0,3.2);
  \draw (3.0,3.2) to[short] (3.0,2.3)
    to[lamp, l={$I_1 = \qty{0.30}{A}$}] (6.2,2.3)
    to[short] (6.2,3.2);
  \draw (3.0,3.2) to[short] (3.0,4.1)
    to[lamp, l={$I_2 = \qty{0.20}{A}$}] (6.2,4.1)
    to[short] (6.2,3.2);
  \draw (6.2,3.2) to[short] (7.8,3.2) to[short] (7.8,0)
    to[short] (0,0);
  \fill (3.0,3.2) circle (2pt);
  \fill (6.2,3.2) circle (2pt);
\end{tikzpicture}

{\small A lei dos \omterm{def:g5:series-parallel-circuits:parallel}{nós} com números: \qty{0.50}{A} chegam à
bifurcação, $0.30 + 0.20$ saem pelos \omterm{def:g5:series-parallel-circuits:parallel}{ramos} --- e \qty{0.50}{A} se
remontam no \omterm{def:g5:series-parallel-circuits:parallel}{nó} lá do outro lado.}
\end{omfigure}

\section{As leis da tensão}

\begin{proposition}[Lei da tensão numa volta em série]\label{prop:g8:circuit-laws:seriesU}
Ao longo de uma volta \omterm{def:g4:series-circuits:series}{em série}, a \omterm{def:g8:voltage-voltmeter:voltage}{tensão} da \omterm{def:g3:first-electric-circuit:battery}{pilha} é repartida entre
os aparelhos, e as parcelas somam exatamente:
\[
U_{\text{pilha}} = U_1 + U_2 + \dots
\]
A \omterm{prop:g5:mirrors-reflection:image}{imagem} da cachoeira faz a contabilidade: a \omterm{def:g3:first-electric-circuit:battery}{pilha} levanta a marcha
pela altura inteira dela, e os aparelhos gastam essa altura entre si,
gota a gota --- com os bons fios não gastando nada.
\end{proposition}

\begin{proposition}[Lei da tensão em ramos paralelos]\label{prop:g8:circuit-laws:parallelU}
Aparelhos \omterm{def:g5:series-parallel-circuits:parallel}{em paralelo} se estendem todos entre os mesmos dois \omterm{def:g5:series-parallel-circuits:parallel}{nós} ---
então uma única \omterm{def:g8:voltage-voltmeter:voltage}{tensão} serve todos eles:
\[
U_1 = U_2 = U_{\text{entre o par de nós}} .
\]
Cada \omterm{def:g5:series-parallel-circuits:parallel}{ramo} desfruta do empurrão inteiro entre os \omterm{def:g5:series-parallel-circuits:parallel}{nós} --- e é por isso
que lâmpadas \omterm{def:g5:series-parallel-circuits:parallel}{em paralelo} brilham no brilho nominal cheio e que cada
tomada da sua casa oferece os mesmos \qty{230}{V}.
\end{proposition}

\begin{omfigure}
\begin{tikzpicture}[scale=0.85]
  \draw (0,0) to[battery1, l=\qty{4.5}{V}] (0,3.0)
    to[short] (1.8,3.0)
    to[lamp, l={$U_1 = \qty{2.1}{V}$}] (4.6,3.0)
    to[lamp, l={$U_2 = \qty{2.4}{V}$}] (7.4,3.0)
    to[short] (8.8,3.0)
    to[short] (8.8,0)
    to[short] (0,0);
\end{tikzpicture}

{\small A lei da \omterm{def:g8:voltage-voltmeter:voltage}{tensão} \omterm{def:g4:series-circuits:series}{em série} com números: os \qty{4.5}{V} da
\omterm{def:g3:first-electric-circuit:battery}{pilha} são gastos em duas parcelas desiguais --- $2.1 + 2.4 = 4.5$,
até o décimo.}
\end{omfigure}

\begin{example}[Por que parcelas desiguais?]\label{ex:g8:circuit-laws:unequal}
Duas lâmpadas diferentes \omterm{def:g4:series-circuits:series}{em série} carregam a mesma corrente (lei um)
e mesmo assim dividem a \omterm{def:g8:voltage-voltmeter:voltage}{tensão} de forma desigual --- \qty{2.1}{V}
contra \qty{2.4}{V}. A parcela maior vai para a lâmpada que se opõe
mais à marcha: quem trabalha mais duro fica com a queda maior. O que
exatamente quer dizer ``se opõe mais'', e como \omterm{def:g6:measurement-in-science:measuring}{medir} isso, é o
assunto inteiro do próximo capítulo --- as leis deste capítulo são o
palco dele.
\end{example}

\section{Calcular antes de medir}

\begin{method}[Resolver um circuito com as quatro leis]\label{met:g8:circuit-laws:solve}
Dado um circuito com algumas leituras conhecidas:
\begin{enumerate}
  \item desenhe o esquema e marque cada $I$ e cada $U$ conhecidos;
  \item escreva as leis aplicáveis como equações --- somas de \omterm{def:g5:series-parallel-circuits:parallel}{nó} para
        as correntes, somas de volta para as tensões, igualdade para
        o $I$ \omterm{def:g4:series-circuits:series}{em série} e para o $U$ \omterm{def:g5:series-parallel-circuits:parallel}{em paralelo};
  \item resolva para as incógnitas (as equações do seu curso de
        álgebra, ganhando o salário delas);
  \item confira o bom senso: lâmpada nenhuma recebendo mais do que a
        \omterm{def:g3:first-electric-circuit:battery}{pilha} oferece, corrente de \omterm{def:g5:series-parallel-circuits:parallel}{ramo} nenhuma passando da corrente
        da estrada principal.
\end{enumerate}
\end{method}

\begin{example}[Resolvido: o circuito misto]\label{ex:g8:circuit-laws:worked}
Uma \omterm{def:g3:first-electric-circuit:battery}{pilha} de \qty{6}{V} alimenta a lâmpada A \omterm{def:g4:series-circuits:series}{em série} com um par \omterm{def:g5:series-parallel-circuits:parallel}{em
paralelo} (B e C). Medidos: $I_{\text{pilha}} = \qty{0.5}{A}$,
$I_B = \qty{0.3}{A}$, $U_A = \qty{2.5}{V}$. Calcule o resto. Lei dos
\omterm{def:g5:series-parallel-circuits:parallel}{nós}: $I_C = 0.5 - 0.3 = \qty{0.2}{A}$. Lei da corrente \omterm{def:g4:series-circuits:series}{em série}: a
lâmpada A carrega os \qty{0.5}{A} inteiros. Lei da \omterm{def:g8:voltage-voltmeter:voltage}{tensão} na volta: o
par recebe $U_{\text{par}} = 6 - 2.5 = \qty{3.5}{V}$; lei da \omterm{def:g8:voltage-voltmeter:voltage}{tensão}
\omterm{def:g5:series-parallel-circuits:parallel}{em paralelo}: $U_B = U_C = \qty{3.5}{V}$. Todas as leituras de
aparelho do circuito, agora conhecidas --- quatro leis, uma linha de
álgebra cada.
\end{example}

\begin{example}[As leis como detectores de mentira]\label{ex:g8:circuit-laws:liedetector}
Um relatório de laboratório afirma: \omterm{def:g3:first-electric-circuit:battery}{pilha} de \qty{4.5}{V}; duas
lâmpadas \omterm{def:g4:series-circuits:series}{em série} com \qty{2.1}{V} e \qty{2.9}{V}. A lei da volta
soma as parcelas: \qty{5.0}{V} gastos de uma bolsa de \qty{4.5}{V}
--- impossível; alguém leu um mostrador errado. Outro relatório: no
\omterm{def:g5:series-parallel-circuits:parallel}{nó} entram \qty{0.40}{A} e pelos \omterm{def:g5:series-parallel-circuits:parallel}{ramos} saem \qty{0.25}{A} e
\qty{0.20}{A} --- a bifurcação ``cria'' \qty{0.05}{A}: rejeitado.
Antes das leis, leituras erradas pareciam tão boas quanto as certas;
agora a constituição audita cada página.
\end{example}

\begin{remark}[Pilhas debaixo das mesmas leis]\label{rem:g8:circuit-laws:batteries}
A regra do uma-atrás-da-outra é a lei da \omterm{def:g8:voltage-voltmeter:voltage}{tensão} \omterm{def:g4:series-circuits:series}{em série} vestindo a
roupa mais velha dela: \omterm{def:g3:first-electric-circuit:battery}{pilhas} \omterm{def:g4:series-circuits:series}{em série} somam as tensões
($1.5 + 1.5 + 1.5 = \qty{4.5}{V}$ --- o segredo da bateria chata), e
uma \omterm{def:g3:first-electric-circuit:battery}{pilha} invertida \emph{subtrai}, com o empurrão dela gasto em
brigar com as outras. As leis governam igualmente quem dá e quem
gasta; só o sinal da contribuição difere.
\end{remark}

\section{Exercícios}

\begin{exercise}[$\star$]\label{exo:g8:circuit-laws:1}
Enuncie as quatro leis --- duas para a corrente, duas para a \omterm{def:g8:voltage-voltmeter:voltage}{tensão}
--- cada uma em uma linha.
\end{exercise}

\begin{exercise}[$\star$]\label{exo:g8:circuit-laws:2}
Um \omterm{def:g5:series-parallel-circuits:parallel}{nó} recebe \qty{0.8}{A} e manda \qty{0.55}{A} por um \omterm{def:g5:series-parallel-circuits:parallel}{ramo}. O outro
\omterm{def:g5:series-parallel-circuits:parallel}{ramo} carrega quanto?
\end{exercise}

\begin{exercise}[$\star$]\label{exo:g8:circuit-laws:3}
Uma bateria de \qty{9}{V} alimenta três lâmpadas \omterm{def:g4:series-circuits:series}{em série}; duas delas
ficam com \qty{2.5}{V} e \qty{3.8}{V}. A parcela da terceira?
\end{exercise}

\begin{exercise}[$\star$]\label{exo:g8:circuit-laws:4}
Duas lâmpadas \omterm{def:g5:series-parallel-circuits:parallel}{em paralelo} numa bateria de \qty{12}{V}: a \omterm{def:g8:voltage-voltmeter:voltage}{tensão} de
cada lâmpada? Que fato do dia a dia sobre as tomadas de casa é a
mesma lei?
\end{exercise}

\begin{exercise}[$\star$]\label{exo:g8:circuit-laws:5}
No circuito misto resolvido, qual medida a mais teria sido
redundante, e por quê? (Escolha qualquer uma e justifique por uma
lei.)
\end{exercise}

\begin{exercise}[$\star$]\label{exo:g8:circuit-laws:6}
Duas lâmpadas diferentes \omterm{def:g4:series-circuits:series}{em série} carregam a mesma corrente mas
dividem a \omterm{def:g8:voltage-voltmeter:voltage}{tensão} em $1.8$ contra \qty{2.7}{V}. Qual lâmpada se opõe
mais à marcha? Que capítulo por vir mede essa oposição?
\end{exercise}

\begin{exercise}[$\star$]\label{exo:g8:circuit-laws:7}
Auditoria: \omterm{def:g3:first-electric-circuit:battery}{pilha} de \qty{6}{V}; lâmpadas \omterm{def:g4:series-circuits:series}{em série} medidas em
\qty{2.2}{V}, \qty{1.9}{V}, \qty{1.4}{V}. As contas fecham dentro da
honestidade de um aparelho?
\end{exercise}

\begin{exercise}[$\star$]\label{exo:g8:circuit-laws:8}
Auditoria: estrada principal \qty{0.9}{A}; três \omterm{def:g5:series-parallel-circuits:parallel}{ramos} com
\qty{0.4}{A}, \qty{0.3}{A}, \qty{0.3}{A}. Veredito?
\end{exercise}

\begin{exercise}[$\star\star$]\label{exo:g8:circuit-laws:9}
Uma bateria de \qty{4.5}{V}, a lâmpada D \omterm{def:g4:series-circuits:series}{em série} com um \omterm{ex:g3:first-electric-circuit:switch}{interruptor}
--- então, \omterm{def:g5:series-parallel-circuits:parallel}{em paralelo} com a lâmpada D, um \omterm{def:g8:voltage-voltmeter:voltmeter}{voltímetro} mostra
\qty{4.5}{V} atravessados no \omterm{ex:g3:first-electric-circuit:switch}{interruptor} \emph{aberto} e \qty{0}{V}
atravessados em D. Mostre que as duas leituras obedecem à lei da
\omterm{def:g8:voltage-voltmeter:voltage}{tensão} na volta.
\end{exercise}

\begin{exercise}[$\star\star$]\label{exo:g8:circuit-laws:10}
Três \omterm{def:g3:first-electric-circuit:battery}{pilhas} iguais uma atrás da outra acendem uma lâmpada; uma
quarta, invertida, entra na fila. \omterm{def:g8:voltage-voltmeter:voltage}{Tensão} total antes e depois, pela
lei da série com sinais? O que a lâmpada faz?
\end{exercise}

\begin{exercise}[$\star\star$]\label{exo:g8:circuit-laws:11}
Um circuito de dois \omterm{def:g5:series-parallel-circuits:parallel}{ramos}: o \omterm{def:g5:series-parallel-circuits:parallel}{ramo} um carrega as lâmpadas E e F \omterm{def:g4:series-circuits:series}{em
série}; o \omterm{def:g5:series-parallel-circuits:parallel}{ramo} dois carrega a lâmpada G sozinha. Bateria de
\qty{6}{V}; $U_E = \qty{2.2}{V}$. Calcule $U_F$ e $U_G$, citando uma
lei para cada passo.
\end{exercise}

\begin{exercise}[$\star\star\star$]\label{exo:g8:circuit-laws:12}
A auditoria completa: uma bateria de \qty{6}{V}; a lâmpada P \omterm{def:g4:series-circuits:series}{em série}
com um \omterm{def:g5:series-parallel-circuits:parallel}{nó}; os \omterm{def:g5:series-parallel-circuits:parallel}{ramos} Q e R se reencontram e voltam. Medidos:
$I_P = \qty{0.6}{A}$, $I_Q = \qty{0.45}{A}$, $U_Q = \qty{2.3}{V}$.
Calcule todas as correntes e tensões restantes do circuito ($I_R$,
$U_P$, $U_R$), enunciando a lei por trás de cada linha --- e então
invente uma ``medida'' a mais que um estudante desatento poderia
relatar e que a sua solução pronta condenaria na hora.
\end{exercise}

\section{Problema: As caixas misteriosas lacradas}

\begin{problem}[{Problema de fim de semana --- o desafio das caixas
lacradas do clube de eletrônica; quatro leis contra quatro mistérios;
a grande auditoria}]\label{pb:g8:circuit-laws:1}
O clube lacra circuitos simples dentro de caixas, deixando de fora
apenas os bornes para os aparelhos. As equipes precisam deduzir o
interior --- constituição na mão. A \omterm{def:g3:first-electric-circuit:battery}{pilha} de bancada é de \qty{6}{V}
o tempo todo.

\medskip\noindent\textbf{Parte I --- Caixa A: dois \omterm{def:g3:first-electric-circuit:battery}{terminais}, duas
lâmpadas.}
A etiqueta admite: duas lâmpadas iguais e nada mais, \omterm{def:g4:series-circuits:series}{em série} ou \omterm{def:g5:series-parallel-circuits:parallel}{em
paralelo}.
\begin{enumerate}
  \item Ligada à \omterm{def:g3:first-electric-circuit:battery}{pilha}, a caixa A puxa \qty{0.2}{A}; sabe-se que uma
        única lâmpada dessas sozinha em \qty{6}{V} puxa cerca de
        \qty{0.4}{A}. Série ou paralelo? Argumente com uma lei da
        corrente.
  \item Preveja o que a caixa puxaria no \emph{outro} arranjo.
  \item Que \omterm{def:g8:voltage-voltmeter:voltage}{tensão} cada lâmpada interna recebe no arranjo real, por
        qual lei da \omterm{def:g8:voltage-voltmeter:voltage}{tensão}?
  \item Em qual arranjo as lâmpadas escondidas brilhariam no brilho
        nominal cheio --- e por que visivelmente não brilham, pelo
        olho mágico da caixa?
\end{enumerate}

\medskip\noindent\textbf{Parte II --- Caixa B: o quebra-cabeça de
três lâmpadas.}
A caixa B admite: uma lâmpada (X) \omterm{def:g4:series-circuits:series}{em série} com um par \omterm{def:g5:series-parallel-circuits:parallel}{em paralelo} (Y,
Z --- não necessariamente iguais).
\begin{enumerate}[resume]
  \item A equipe mede $I_{\text{pilha}} = \qty{0.5}{A}$ e, no borne
        de teste da própria Y, $I_Y = \qty{0.35}{A}$. Ache $I_Z$ e
        $I_X$, citando leis.
  \item Um \omterm{def:g8:voltage-voltmeter:voltmeter}{voltímetro} atravessado em X: \qty{2.8}{V}. Ache a \omterm{def:g8:voltage-voltmeter:voltage}{tensão}
        do par \omterm{def:g5:series-parallel-circuits:parallel}{em paralelo} --- e cada uma de $U_Y$ e $U_Z$.
  \item Y e Z carregam correntes diferentes com a mesma \omterm{def:g8:voltage-voltmeter:voltage}{tensão}. O que
        isso revela sobre as duas lâmpadas, na linguagem do
        \cref{ex:g8:circuit-laws:unequal}?
  \item A equipe desenrosca Y por uma portinhola. Preveja o sentido
        da mudança da nova $I_{\text{pilha}}$ (sobe, cai ou fica),
        raciocinando com estradas e a oposição delas.
\end{enumerate}

\medskip\noindent\textbf{Parte III --- Caixa C: a trapaça.}
\begin{enumerate}[resume]
  \item O construtor da caixa C relata: ``\omterm{def:g3:first-electric-circuit:battery}{pilha} de \qty{6}{V}; duas
        lâmpadas \omterm{def:g4:series-circuits:series}{em série} lá dentro com \qty{2.5}{V} e \qty{4.1}{V};
        \omterm{def:g5:series-parallel-circuits:parallel}{nó} interno recebendo \qty{0.3}{A} e mandando \qty{0.2}{A}
        mais \qty{0.15}{A}.'' Condene o relatório duas vezes, uma lei
        por condenação.
  \item Uma das duas leituras de \omterm{def:g8:voltage-voltmeter:voltage}{tensão} relatadas se revela depois
        honesta. Se for a de \qty{2.5}{V}, quanto a parceira dela
        precisa valer de verdade?
  \item O construtor confessa que os números do \omterm{def:g5:series-parallel-circuits:parallel}{nó} foram
        ``arredondados com generosidade''. Se as duas leituras dos
        \omterm{def:g5:series-parallel-circuits:parallel}{ramos}, \qty{0.2}{A} e \qty{0.15}{A}, forem honestas, quanto a
        estrada principal precisa carregar de verdade?
  \item Por que as quatro leis fazem da trapaça de caixa lacrada um
        jogo perdido? Uma frase.
\end{enumerate}

\medskip\noindent\textbf{Parte IV --- A constituição.}
\begin{enumerate}[resume]
  \item Escreva a carta do clube: as quatro leis em quatro linhas,
        cada uma com a \omterm{prop:g5:mirrors-reflection:image}{imagem} do dia a dia que a carrega (estradas,
        pedágios, cachoeiras, bifurcações).
\end{enumerate}
\end{problem}
